Pour chacune des affirmations suivantes, indiquer si elle est vraie ou fausse. Justifier la réponse donnée. Une réponse non justifiée ne rapporte aucun point.

L'espace est muni d'un repère orthonormé \Oijk.

On considère~:

\begin{itemize}[noitemsep]
\item les points $S(-1~;~\sqrt{2}~;-4)$, A$(2~;~\sqrt{2}~;~-1)$, B$(1~;~\sqrt{2}~;0)$ et C$(2~;~0~;-1)$~;
\item la droite $d$ dont une représentation paramétrique est donnée par~:
\[
\begin{cases}
x = 2-k\\
y = \sqrt{2}\,k\\
z = -1-k
\end{cases}, \quad k \in \R~;
\]
\item le plan $P$ dont une équation cartésienne est~: $x - z + 1 = 0$.
\end{itemize}

On admet que les points A, B et C ne sont pas alignés.

\textbf{Affirmation 1~:} Le vecteur $\overrightarrow{\text{SA}}$ est un vecteur normal au plan (ABC).

\textbf{Affirmation 2~:} Les droites (SB) et $d$ sont sécantes.

\textbf{Affirmation 3~:} La droite $d$ est parallèle au plan $P$.

\textbf{Affirmation 4~:} Le projeté orthogonal de S sur le plan $P$ est le point H de coordonnées $(-3~;~\sqrt{2}~;~-2)$.
